\documentclass[aspectratio=169,10pt]{beamer}
\usetheme[numbering=fraction,progressbar=frametitle]{metropolis}
\usepackage{fontspec}
\setsansfont{Latin Modern Sans}
\usepackage{amsmath,amssymb,mathtools,graphicx,booktabs}
\usepackage{polyglossia}
\setdefaultlanguage{french}
\setbeamertemplate{navigation symbols}{}
\setbeamerfont{frametitle}{size=\large}
\setbeamerfont{normal text}{size=\normalsize}
\setlength{\parskip}{.45em}
\hypersetup{pdfauthor={},pdfsubject={Lecture slides},pdfkeywords={},pdftitle={Divisibilité, nombres premiers, PGCD et PPCM}}
\providecommand{\N}{\mathbb{N}}
\providecommand{\Z}{\mathbb{Z}}
\providecommand{\D}{\mathbb{D}}
\providecommand{\Q}{\mathbb{Q}}
\providecommand{\R}{\mathbb{R}}
\providecommand{\C}{\mathbb{C}}
\providecommand{\vect}[1]{\overrightarrow{#1}}
\providecommand{\vi}{\vec{\imath}}
\providecommand{\vj}{\vec{\jmath}}
\providecommand{\vk}{\vec{k}}
\providecommand{\norm}[1]{\left\lVert #1\right\rVert}
\providecommand{\vc}[2]{\begin{pmatrix}#1\\#2\end{pmatrix}}
\providecommand{\deter}[2]{\det\left(#1,#2\right)}
\providecommand{\abs}[1]{\left\lvert #1\right\rvert}
\providecommand{\dd}{\,\mathrm{d}}
\providecommand{\e}{\mathrm{e}}
\providecommand{\Card}{\operatorname{Card}}
\providecommand{\E}{\operatorname{E}}
\providecommand{\V}{\operatorname{V}}
\providecommand{\degre}{^{\circ}}
\providecommand{\un}[1]{\,\mathrm{#1}}
\providecommand{\sci}[2]{#1\times10^{#2}}
\providecommand{\pgcd}{\operatorname{PGCD}}
\providecommand{\ppcm}{\operatorname{PPCM}}
\title{Divisibilité, nombres premiers, PGCD et PPCM}
\subtitle{Cours et exemples}
\date{}
\author{}
\begin{document}
\begin{frame}[plain]
\vfill
{\usebeamerfont{title}\inserttitle\par}
\medskip
{\usebeamerfont{subtitle}\insertsubtitle\par}
\medskip
{\usebeamercolor[fg]{progress bar}\rule{\textwidth}{.5pt}}
\vfill
\end{frame}
\begin{frame}{Objectifs}
Décomposer un entier en facteurs premiers, utiliser Euclide et résoudre des problèmes de regroupement et de périodicité.
\end{frame}
\begin{frame}{Diviseurs et nombres premiers}
Pour \(a,b\in\mathbb{Z}\), \(a\mid b\) signifie qu’il existe \(k\in\mathbb{Z}\) avec \(b=ak\). Un entier naturel \(p\ge2\) est premier s’il a exactement deux diviseurs positifs : \(1\) et \(p\). Le nombre \(1\) n’est pas premier.
\end{frame}
\begin{frame}{Diviseurs et nombres premiers}
Si \(a\mid b\) et \(a\mid c\), alors \(a\mid (ub+vc)\) pour tous entiers \(u,v\). Si \(n\ge2\) est composé, il possède un diviseur premier inférieur ou égal à \(\sqrt n\).
\end{frame}
\begin{frame}{Diviseurs et nombres premiers}
Écrire \(b=ak,c=a\ell\) donne \(ub+vc=a(uk+v\ell)\). Si \(n=de\) avec \(1<d\le e<n\), alors \(d^2\le n\); \(d\) possède un diviseur premier (en prenant son plus petit diviseur supérieur à \(1\)), lui-même au plus \(\sqrt n\).
\end{frame}
\begin{frame}{Diviseurs et nombres premiers}
Tout entier \(n\ge2\) possède une décomposition en produit de nombres premiers, unique à l’ordre des facteurs près.
\end{frame}
\begin{frame}{Diviseurs et nombres premiers}
\(360=2^3\times3^2\times5\). Pour tester \(97\), il suffit de vérifier les nombres premiers \(2,3,5,7\), car \(\sqrt{97}<10\). Aucun ne divise \(97\), donc \(97\) est premier.
\end{frame}
\begin{frame}{PGCD et algorithme d’Euclide}
Pour \(a,b\in\mathbb{N}^*\), \(\operatorname{PGCD}(a,b)\) est le plus grand diviseur positif commun. Ils sont premiers entre eux lorsque ce PGCD vaut \(1\). On étend par \(\operatorname{PGCD}(a,0)=\operatorname{PGCD}(0,a)=a\) pour \(a>0\); \(\operatorname{PGCD}(0,0)\) n’est pas utilisé ici.
\end{frame}
\begin{frame}{PGCD et algorithme d’Euclide}
Si \(a=bq+r\), les diviseurs communs de \(a,b\) sont ceux de \(b,r\). Ainsi \(\operatorname{PGCD}(a,b)=\operatorname{PGCD}(b,r)\).
\end{frame}
\begin{frame}{PGCD et algorithme d’Euclide}
Un diviseur de \(a,b\) divise \(r=a-bq\). Un diviseur de \(b,r\) divise \(a=bq+r\). Les ensembles sont les mêmes.
\end{frame}
\begin{frame}{PGCD et algorithme d’Euclide}
\(252=198+54\); \(198=3\times54+36\); \(54=36+18\); \(36=2\times18\). Le dernier reste non nul est \(18\), donc \(\operatorname{PGCD}(252,198)=18\).
\end{frame}
\begin{frame}{PGCD et algorithme d’Euclide}
Entrées \(a,b\in\mathbb{N}\), non tous deux nuls. Tant que \(b\ne0\), calculer \(r\leftarrow a\bmod b\), puis \(a\leftarrow b\), \(b\leftarrow r\). Afficher \(a\). Les restes non nuls décroissent strictement, donc la boucle termine. Les diviseurs communs sont conservés; la sortie est le PGCD.
\end{frame}
\begin{frame}{PPCM}
Pour \(a,b\in\mathbb{N}^*\), \(\operatorname{PPCM}(a,b)\) est le plus petit multiple positif commun.
\end{frame}
\begin{frame}{PPCM}
Dans les décompositions premières, le PGCD prend les exposants minimaux et le PPCM les exposants maximaux (un facteur absent a exposant \(0\)). Ainsi \(\operatorname{PGCD}(a,b)\operatorname{PPCM}(a,b)=ab\).
\end{frame}
\begin{frame}{PPCM}
Un diviseur commun ne peut avoir un exposant supérieur au minimum; le produit des minima convient. Un multiple commun doit avoir au moins les maxima; leur produit est le plus petit. Pour chaque premier, minimum plus maximum égale la somme des deux exposants.
\end{frame}
\begin{frame}{PPCM}
\(72=2^3 3^2\), \(90=2\,3^2 5\) : PGCD \(18\), PPCM \(360\). Pour trois entiers, calculer successivement le PGCD ou PPCM de deux, puis avec le troisième.
\end{frame}
\begin{frame}{PPCM}
Pour des lots identiques sans reste, le nombre de lots doit diviser chaque effectif; le plus grand nombre demande donc un PGCD. Pour un retour commun de cycles partant ensemble, la durée doit être un multiple de chaque période; le premier retour positif demande un PPCM. Justifier cette divisibilité avant d’appliquer l’algorithme.
\end{frame}
\begin{frame}{PPCM}
Avec \(48\) objets rouges et \(60\) bleus, \(\operatorname{PGCD}(48,60)=12\) donne au maximum \(12\) lots, de \(4\) rouges et \(5\) bleus chacun. Pour des cycles de \(48\) et \(60\) minutes démarrant ensemble, \(\operatorname{PPCM}(48,60)=240\) donne le premier retour commun après \(240\) minutes. Intervertir les deux nombres aurait donné des réponses sans rapport avec les contraintes.
\end{frame}
\begin{frame}{PPCM}
Pour \(n\ge2\), chercher les diviseurs premiers jusqu’à \(\sqrt n\). Si un diviseur est trouvé, donner une factorisation non triviale. Si aucun n’est trouvé, expliquer pourquoi un éventuel facteur composé aurait lui-même un facteur premier dans cette plage. Arrêter les tests à une borne sans la justifier ne prouve pas la primalité.
\end{frame}
\begin{frame}{PPCM}
Pour \(149\), \(12^2<149<13^2\). Les seuls premiers à tester sont \(2,3,5,7,11\). Les trois premiers sont exclus par les critères usuels; \(149=7\times21+2\) et \(149=11\times13+6\). Aucun ne divise \(149\) : ce nombre est premier.
\end{frame}
\end{document}
